\section{Erd\H{o}s Problem \#151: independent continuation}
\noindent Date: 23 September 2026. Source versions read in full: \texttt{151.tex} in \texttt{gpt\_pro\_5.4}.

\subsection*{1) OBJECTIVE AND SOURCE AUDIT}
Let $\tau(G)$ hit every maximal clique with at least two vertices, and let
$H(n)$ be the minimum independence number among triangle-free $n$-vertex
graphs. The target is $\tau(G)\le n-H(n)$.
The source codegree bound is correct, but its proof only uses one endpoint:
the stronger bound $\tau(G)\le n-\Delta(G)$ follows immediately. The asserted
new codegree obstruction is therefore redundant with this simpler bound.

\subsection*{2) WORK: A LOCAL-PLUS-REMOTE CERTIFICATE}
For a vertex $v$ put $U=N(v)$, and let $W$ consist of vertices outside
$U\cup\{v\}$ which have no neighbour in $U$. Then
\[
 \tau(G)\le n-|U|-\alpha(G[W]). \tag{1}
\]
Take an independent set $I\subset W$ of size $\alpha(G[W])$ and set
$T=V(G)\setminus(U\cup I)$. There are no edges between $U$ and $I$ and
no edges inside $I$. Consequently a clique of size at least two contained
in $U\cup I$ must lie entirely in $U$. It extends by adding $v$, so is
not maximal in $G$. Thus $T$ is a transversal, proving (1).
The convention excluding singleton cliques is important here.

A computable weaker form is
\[
 \tau(G)\le n-d(v)-\frac{|W|}{\Delta(G[W])+1},
\]
with the last term zero when $W$ is empty. Indeed, greedily select a vertex
of $G[W]$ and remove it and its neighbours; at most $\Delta(G[W])+1$
vertices are removed at each step, leaving an independent set of the
claimed size (with an integer ceiling if desired).

Hence a counterexample must satisfy, for every $v$,
\[
 d(v)+\alpha(G[W_v])<H(n),
\]
as well as $\alpha(G)<H(n)$. This strengthens the supplied local certificate
by adding a compatible remote independent set. It is a sufficient condition,
not a characterization of all clique transversals.

\subsection*{3) ADVERSARIAL VERIFICATION}
One cannot add an arbitrary independent set outside $N(v)$: it may have
edges into $N(v)$ and create a maximal clique entirely inside the proposed
avoided set. The definition of $W$ is exactly what prevents this failure.
In a complete graph (1) gives $\tau\le1$; with an isolated $v$ it reduces
to an independence-based bound. Neither boundary case creates singleton
clique obligations excluded by the problem.

\subsection*{7) FINAL AND REFERENCE}
\textbf{UNRESOLVED.} The certificate improves on the source codegree argument,
but I cannot prove its left side always reaches $H(n)$, or construct a
graph violating the full conjecture. The definition of $H(n)$ and the
clique convention were checked at
\url{https://www.erdosproblems.com/latex/151}.

\subsection*{8) COMPLETION ESTIMATE}
\textbf{COMPLETION: 10\%}
This is a subjective estimate of progress toward the entire stated problem, not a probability that the conjecture or an unverified proof is correct.
